\documentclass[a4paper,12pt,twoside]{memoir}

% Castellano
\usepackage[spanish,es-tabla]{babel}
\selectlanguage{spanish}
\usepackage{graphicx}
\usepackage{float}
\usepackage[utf8]{inputenc}
\usepackage[T1]{fontenc}
\usepackage{lmodern} % scalable font
\usepackage{microtype}
\usepackage{placeins}

\RequirePackage{booktabs}
\RequirePackage[table]{xcolor}
\RequirePackage{xtab}
\RequirePackage{multirow}

% Links
\PassOptionsToPackage{hyphens}{url}\usepackage[colorlinks]{hyperref}
\hypersetup{
	allcolors = {red}
}

% Ecuaciones
\usepackage{amsmath}

% Rutas de fichero / paquete
\newcommand{\ruta}[1]{{\sffamily #1}}

% Párrafos
\nonzeroparskip

% Huérfanas y viudas
\widowpenalty100000
\clubpenalty100000

% Evitar solapes en el header
\nouppercaseheads

% Imagenes
\usepackage{graphicx}
\newcommand{\imagen}[2]{
	\begin{figure}[!h]
		\centering
		\includegraphics[width=0.9\textwidth]{#1}
		\caption{#2}\label{fig:#1}
	\end{figure}
	\FloatBarrier
}

\newcommand{\imagenflotante}[2]{
	\begin{figure}%[!h]
		\centering
		\includegraphics[width=0.9\textwidth]{#1}
		\caption{#2}\label{fig:#1}
	\end{figure}
}



% El comando \figura nos permite insertar figuras comodamente, y utilizando
% siempre el mismo formato. Los parametros son:
% 1 -> Porcentaje del ancho de página que ocupará la figura (de 0 a 1)
% 2 --> Fichero de la imagen
% 3 --> Texto a pie de imagen
% 4 --> Etiqueta (label) para referencias
% 5 --> Opciones que queramos pasarle al \includegraphics
% 6 --> Opciones de posicionamiento a pasarle a \begin{figure}
\newcommand{\figuraConPosicion}[6]{%
  \setlength{\anchoFloat}{#1\textwidth}%
  \addtolength{\anchoFloat}{-4\fboxsep}%
  \setlength{\anchoFigura}{\anchoFloat}%
  \begin{figure}[#6]
    \begin{center}%
      \Ovalbox{%
        \begin{minipage}{\anchoFloat}%
          \begin{center}%
            \includegraphics[width=\anchoFigura,#5]{#2}%
            \caption{#3}%
            \label{#4}%
          \end{center}%
        \end{minipage}
      }%
    \end{center}%
  \end{figure}%
}

%
% Comando para incluir imágenes en formato apaisado (sin marco).
\newcommand{\figuraApaisadaSinMarco}[5]{%
  \begin{figure}%
    \begin{center}%
    \includegraphics[angle=90,height=#1\textheight,#5]{#2}%
    \caption{#3}%
    \label{#4}%
    \end{center}%
  \end{figure}%
}
% Para las tablas
\newcommand{\otoprule}{\midrule [\heavyrulewidth]}
%
% Nuevo comando para tablas pequeñas (menos de una página).
\newcommand{\tablaSmall}[5]{%
 \begin{table}
  \begin{center}
   \rowcolors {2}{gray!35}{}
   \begin{tabular}{#2}
    \toprule
    #4
    \otoprule
    #5
    \bottomrule
   \end{tabular}
   \caption{#1}
   \label{tabla:#3}
  \end{center}
 \end{table}
}

%
%Para el float H de tablaSmallSinColores
\usepackage{float}

%
% Nuevo comando para tablas pequeñas (menos de una página).
\newcommand{\tablaSmallSinColores}[5]{%
 \begin{table}[H]
  \begin{center}
   \begin{tabular}{#2}
    \toprule
    #4
    \otoprule
    #5
    \bottomrule
   \end{tabular}
   \caption{#1}
   \label{tabla:#3}
  \end{center}
 \end{table}
}

\newcommand{\tablaApaisadaSmall}[5]{%
\begin{landscape}
  \begin{table}
   \begin{center}
    \rowcolors {2}{gray!35}{}
    \begin{tabular}{#2}
     \toprule
     #4
     \otoprule
     #5
     \bottomrule
    \end{tabular}
    \caption{#1}
    \label{tabla:#3}
   \end{center}
  \end{table}
\end{landscape}
}

%
% Nuevo comando para tablas grandes con cabecera y filas alternas coloreadas en gris.
\newcommand{\tabla}[6]{%
  \begin{center}
    \tablefirsthead{
      \toprule
      #5
      \otoprule
    }
    \tablehead{
      \multicolumn{#3}{l}{\small\sl continúa desde la página anterior}\\
      \toprule
      #5
      \otoprule
    }
    \tabletail{
      \hline
      \multicolumn{#3}{r}{\small\sl continúa en la página siguiente}\\
    }
    \tablelasttail{
      \hline
    }
    \bottomcaption{#1}
    \rowcolors {2}{gray!35}{}
    \begin{xtabular}{#2}
      #6
      \bottomrule
    \end{xtabular}
    \label{tabla:#4}
  \end{center}
}

%
% Nuevo comando para tablas grandes con cabecera.
\newcommand{\tablaSinColores}[6]{%
  \begin{center}
    \tablefirsthead{
      \toprule
      #5
      \otoprule
    }
    \tablehead{
      \multicolumn{#3}{l}{\small\sl continúa desde la página anterior}\\
      \toprule
      #5
      \otoprule
    }
    \tabletail{
      \hline
      \multicolumn{#3}{r}{\small\sl continúa en la página siguiente}\\
    }
    \tablelasttail{
      \hline
    }
    \bottomcaption{#1}
    \begin{xtabular}{#2}
      #6
      \bottomrule
    \end{xtabular}
    \label{tabla:#4}
  \end{center}
}

%
% Nuevo comando para tablas grandes sin cabecera.
\newcommand{\tablaSinCabecera}[5]{%
  \begin{center}
    \tablefirsthead{
      \toprule
    }
    \tablehead{
      \multicolumn{#3}{l}{\small\sl continúa desde la página anterior}\\
      \hline
    }
    \tabletail{
      \hline
      \multicolumn{#3}{r}{\small\sl continúa en la página siguiente}\\
    }
    \tablelasttail{
      \hline
    }
    \bottomcaption{#1}
  \begin{xtabular}{#2}
    #5
   \bottomrule
  \end{xtabular}
  \label{tabla:#4}
  \end{center}
}



\definecolor{cgoLight}{HTML}{EEEEEE}
\definecolor{cgoExtralight}{HTML}{FFFFFF}

%
% Nuevo comando para tablas grandes sin cabecera.
\newcommand{\tablaSinCabeceraConBandas}[5]{%
  \begin{center}
    \tablefirsthead{
      \toprule
    }
    \tablehead{
      \multicolumn{#3}{l}{\small\sl continúa desde la página anterior}\\
      \hline
    }
    \tabletail{
      \hline
      \multicolumn{#3}{r}{\small\sl continúa en la página siguiente}\\
    }
    \tablelasttail{
      \hline
    }
    \bottomcaption{#1}
    \rowcolors[]{1}{cgoExtralight}{cgoLight}

  \begin{xtabular}{#2}
    #5
   \bottomrule
  \end{xtabular}
  \label{tabla:#4}
  \end{center}
}




\graphicspath{ {./img/} }

% Capítulos
\chapterstyle{bianchi}
\newcommand{\capitulo}[2]{
	\setcounter{chapter}{#1}
	\setcounter{section}{0}
	\setcounter{figure}{0}
	\setcounter{table}{0}
	\chapter*{#2}
	\addcontentsline{toc}{chapter}{#2}
	\markboth{#2}{#2}
}

% Apéndices
\renewcommand{\appendixname}{Apéndice}
\renewcommand*\cftappendixname{\appendixname}

\newcommand{\apendice}[1]{
	%\renewcommand{\thechapter}{A}
	\chapter{#1}
}

\renewcommand*\cftappendixname{\appendixname\ }

% Formato de portada
\makeatletter
\usepackage{xcolor}
\newcommand{\tutor}[1]{\def\@tutor{#1}}
\newcommand{\course}[1]{\def\@course{#1}}
\definecolor{cpardoBox}{HTML}{E6E6FF}
\def\maketitle{
  \null
  \thispagestyle{empty}
  % Cabecera ----------------
\noindent\includegraphics[width=\textwidth]{cabecera}\vspace{1cm}%
  \vfill
  % Título proyecto y escudo informática ----------------
  \colorbox{cpardoBox}{%
    \begin{minipage}{.8\textwidth}
      \vspace{.5cm}\Large
      \begin{center}
      \textbf{TFG del Grado en Ingeniería Informática}\vspace{.6cm}\\
      \textbf{\LARGE\@title{}}
      \end{center}
      \vspace{.2cm}
    \end{minipage}

  }%
  \hfill\begin{minipage}{.20\textwidth}
    \includegraphics[width=\textwidth]{escudoInfor}
  \end{minipage}
  \vfill
  % Datos de alumno, curso y tutores ------------------
  \begin{center}%
  {%
    \noindent\LARGE
    Presentado por \@author{}\\ 
    en Universidad de Burgos --- \@date{}\\
    Tutor: \@tutor{}\\
  }%
  \end{center}%
  \null
  \cleardoublepage
  }
\makeatother


% Datos de portada
\title{eLearningQA \\Documentación Técnica}
\author{Roberto Arasti Blanco}
\tutor{Raúl Marticorena Sánchez y Carlos López Nozal}
\date{\today}

\begin{document}

\maketitle



\cleardoublepage



%%%%%%%%%%%%%%%%%%%%%%%%%%%%%%%%%%%%%%%%%%%%%%%%%%%%%%%%%%%%%%%%%%%%%%%%%%%%%%%%%%%%%%%%



\frontmatter


\clearpage

% Indices
\tableofcontents

\clearpage

\listoffigures

\clearpage

\listoftables

\clearpage

\mainmatter

\appendix

\apendice{Plan de Proyecto Software}

\section{Introducción}
En este apéndice se trata sobre los costes y viabilidad del proyecto. Contiene una planificación temporal que hablará de la cantidad de tiempo estimada que va a costar realizar el proyecto. También un estudio de viabilidad dividido en dos partes: un estudio de viabilidad económica que tratará los costes y beneficios teóricos de la realización del proyecto en un escenario real a pesar de que el fin de este trabajo sea educativo, y un estudio de viabilidad legal que discutirá si hay o no problemas derivados de las licencias de las herramientas usadas y/o infracciones de copyright.
\section{Planificación temporal}
En esta sección se habla de la cantidad de tiempo estimada necesaria al principio del proyecto y la cantidad de tiempo que ha acabado llevando, detallando lo ocurrido durante el proceso.

La planificación del proyecto se ha llevado a cabo en sprints de una duración de dos semanas. Normalmente se realizó una reunión entre sprint y sprint en la que se discutían los resultados del sprint anterior y los objetivos para el siguiente. En GitHub, los sprints están definidos como milestones (hitos) y los objetivos del sprint como issues (problemas, asuntos). Para cada sprint he creado un nuevo milestone y he incluido los issues respectivos a los objetivos del sprint a dicho milestone.

Al comienzo estimé que para completar el proyecto necesitaría 400 horas de trabajo a lo largo de 16 semanas, lo que conllevaría 25 horas de trabajo a la semana.

Ahora paso a resumir los objetivos y sucesos de cada sprint:

\subsection{Sprint 0 (7/10/2021 - 21/10/2021)}
	En este sprint los objetivos propuestos eran instalar la plantilla de la memoria en LaTeX proporcionada por la UBU, redactar una introducción y definir los objetivos del proyecto, y sintetizar una lista de comprobaciones a realizar por parte de la aplicación a partir de unos marcos de calidad. En la reunión de cierre del sprint los tutores se mostraron disconformes con mi forma de redactar excesivamente técnica y breve y se propuso el realizar reuniones a mitad de sprint cuando yo lo viera conveniente.
\subsection{Sprint 1 (21/10/2021 - 4/11/2021)}
	En este sprint se definieron como objetivos incluir en la introducción un resumen de la situación actual respecto del control de calidad en el e-learning y añadir las dimensiones de roles y perspectivas a la lista de comprobaciones y añadir la lista a la memoria. En la reunión de mitad de sprint se remarcó la falta de referencias bibliográficas y se me pidió que expandiera el apartado de conceptos teóricos. En la reunión de cierre del sprint los tutores se mostraron bastante más conformes con mi trabajo.
\subsection{Sprint 2 (4/11/2021 - 18/11/2021)}
	En este sprint se definieron como objetivos añadir definiciones al apartado de conceptos teóricos, mostrar la lista de consultas definida en forma de tabla, y decantarme por alguna de los tres frameworks para el desarrollo de la aplicación (Vaadin, JavaFX, o Spring). En la reunión de mitad de sprint se me dieron indicaciones sobre aspectos a mejorar que realicé antes de acabar el sprint. Elegí Spring para crear la aplicación web.
\subsection{Sprint 3 (18/11/2021 - 2/12/2021)}
	En este sprint el objetivo principal fue la creación de un prototipo. Creé un prototipo muy simple que mostraba la lista de cursos recientemente accedidos en la plataforma de demostración Mount Orange School tras introducir usuario y contraseña. En la reunión de final de sprint se me propuso utilizar el patrón de diseño fachada para separar la responsabilidad de acceder a los servicios web y generar el contenido para mostrarlo en HTML.
\subsection{Sprint 4 (2/12/2021 - 16/12/2021)}
	En este sprint se definieron como objetivos comenzar con el diseño de la aplicación, crear un pipeline de integración continua/despliegue continuo con GitHub Actions y Heroku, y añadir funcionalidad al prototipo. En la primera semana del sprint creé un boceto de la interfaz de la aplicación y un diagrama de clases preliminar, además de implementar la integración continua con GitHub Actions y despliegue continuo en Heroku. En la reunión de mitad de sprint mostré los avances que había conseguido a Raúl Marticorena (Carlos López estaba ocupado ese día) y me instó a añadir funciones al prototipo. Durante la segunda semana añadí una versión reducida del informe específico que mostraba la comprobación de si un curso tenía grupos definidos o no. En la reunión de final de sprint mostré los avances que había conseguido a Carlos López (Raúl Marticorena estaba ocupado ese día) y establecimos que no íbamos a poder reunirnos hasta la vuelta de las vacaciones y me motivó diciendo que a partir de ese punto el resultado del proyecto dependía de la cantidad de trabajo que invirtiera en él y que gracias al trabajo que había realizado hasta ese punto el proyecto estaba encarrilado.
\subsection{Sprint 5 (16/12/2021 - 30/12/2021)}
	En este sprint se definieron como objetivos refactorizar la fachada de forma que se puedan realizar tests sobre ella, añadir tests, y seguir implementando consultas sobre los cursos Moodle en el prototipo. Al final de este sprint no hubo reunión por motivo de las fiestas navideñas.
\subsection{Sprint 6 (30/12/2021 - 13/1/2022)}
	En este sprint decidí implementar las comprobaciones restantes, rellenar algunos apartados de la memoria y cambiar el estilo de la página del informe para que hubiera un código de cinco colores para distintos niveles de gravedad. En la reunión, a pesar de haber hecho un gran avance, decidí retrasar la entrega de este trabajo al segundo semestre.
\subsection{Sprint 7 (13/1/2022 - 27/1/2022)}
	En este sprint refactoricé los test unitarios para que carguen los objetos necesarios desde una serie de archivos json en vez de conectarse a la página de demo de Moodle que es accesible por cualquiera y está sujeta a cambios. También se refactorizó la fachada para no tener métodos extremadamente largos y cambios varios. A la reunión de final de sprint solo pudo acudir Carlos López, y propuso añadir cinco consultas que comprobaban aspectos de los cuestionarios.
\subsection{Sprint 8 (27/1/2022 - 10/2/2022)}
	En este sprint integré la herramienta de calidad de código SonarCloud al repositorio, revisé la API de servicios web de Moodle para ver si eran factibles las consultas propuestas en la reunión anterior, creé un archivo html a modo de boceto para trabajar en el diseño del informe e hice un apartado de buenas prácticas en la docencia online. En la reunión mostré el boceto html, con el que se mostraron medianamente conformes. En la reunión también se desestimó el implementar las consultas propuestas y dejarlo como una línea futura de trabajo.
\subsection{Sprint 9 (10/2/2022 - 24/2/2022)}
	En este sprint empecé a solucionar los errores descubiertos por SonarCloud y modifiqué el estilo de la mayoría de mis páginas con el framework Bootstrap. En la reunión, hablando del logo acacabamos descubriendo que al ser ``Moodle'' una marca registrada debíamos hacer cambios como por ejemplo el nombre de dominio del despliegue en Heroku.
\subsection{Sprint 10 (24/2/2022 - 10/3/2022)} 
	En este sprint cambié con Bootstrap el estilo de la página de informe, cambié todos los nombres problemáticos y solucioné casi todos los errores de código. En la reunión se decidió seguir trabajando en el estilo del informe e implementar una forma de generar alertas para los errores detectados por las consultas.
\subsection{Sprint 11 (10/3/2022 - 24/3/2022)} 
	En este sprint añadí las alertas y diseñé el logo de la aplicación. En la reunión se comentaron errores del informe y se decidió añadir gráficos a la página de informe y hacer perfiles de configuración intercambiables para distintos tipos de curso.
\subsection{Sprint 12 (24/3/2022 - 7/4/2022)} 
	En este sprint arreglé algunas de las alertas que no funcionaban bien y añadí los perfiles de configuración, pero aplacé añadir los gráficos.
\subsection{Sprint 13 (7/4/2022 - 21/4/2022)} 
	En este sprint arregle un error encontrado por los profesores, añadí los gráficos y una página de error, pero el cambio necesario para hacer las comprobaciones flexibles tenía demasiada complejidad. Solo Carlos López pudo acudir a la reunión, así que se decidió convocar otra la semana siguiente. 
\subsection{Sprint 14 (21/4/2022 - 28/4/2022)} 
	En este sprint se arreglaron problemas estéticos y pequeños bugs. En la reunión dimos un repaso a todos los errores y mejoras pendientes de la página de informe. También se decidió empezar a centrarse en la memoria y anexos.
\subsection{Sprint 15 (28/4/2022 - 12/5/2022)} 
	En este sprint se arreglaron más errores y se hicieron correcciones en memoria y anexos. En la reunión se comentaron más errores y se decidió continuar con las mismas tareas.
\subsection{Sprint 16 (12/5/2022 - 26/5/2022)} 
	En este sprint se siguió haciendo lo mismo que en el sprint anterior, arreglar errores y añadir contenido a la memoria. En la reunión se decidió pasar a hacer sprints de una semana debido a la cercanía de la fecha limite y centrarse en completar los anexos.
\subsection{Sprint 17 (26/5/2022 - 2/6/2022)} 
	En este sprint se solucionó alguno de los problemas restantes de la aplicación y se rellenaron algunos apartados restantes de los anexos. A la reunión solo pudo acudir Carlos López, y en esta se decidió dejar de arreglar errores de código y centrarse en treminar memoria y anexos.
\subsection{Sprint 18 (2/6/2022 - 9/6/2022)} 
	En este sprint se revisó la memoria a partir de los comentarios recibidos de los tutores. No hubo reunión porque se había decidido comunicarse mediante correos a partir de este punto.

\section{Estudio de viabilidad}

\subsection{Viabilidad económica}
En este apartado se estiman los costes que tendría la realización del proyecto en un caso real.

A pesar de haber llevado 8 meses el desarrollo del proyecto, la cantidad de trabajo ronda las 800 horas, que equivaldría a 5 meses a jornada completa que es en lo que nos basaremos para el cálculo.
\begin{table}[H]
	\centering
	\caption{Costes de personal}
	\resizebox{0.5\textwidth}{!}{%
		\begin{tabular}{lc}
			\hline
			\textbf{Concepto}        & \textbf{Coste} \\ \hline
			Sueldo mensual neto      & 1.000,00€      \\
			Retención IRPF (19\%)    & 254,52€        \\
			Seguridad Social (6,35\%) & 85,06€         \\
			Sueldo mensual bruto     & 1.339,58€      \\ \hline
			\textbf{Total 5 meses}   & 6.697,90€      \\ \hline
		\end{tabular}%
	}
\end{table}
Las cuotas a la seguridad social se componen de un 4,70\% de contingencias comunes, un 1,55\% por desempleo de tipo general, y un 0,10\% de formación profesional.

Se supone que el ordenador portátil se amortiza en cinco años y ha sido usado cinco meses.
\begin{table}[H]
	\centering
	\caption{Costes de hardware}
	\resizebox{0.75\textwidth}{!}{%
		\begin{tabular}{lcc}
			\hline
			\textbf{Concepto}        & \textbf{Coste}& \textbf{Coste amortizado} \\ \hline
			Ordenador portátil     & 550€ & 45,83€    \\ \hline
			\textbf{Total}   & 550€ & 45,83€    \\ \hline
		\end{tabular}%
	}
\end{table}

\begin{table}[H]
	\centering
	\caption{Costes varios}
	\resizebox{0.5\textwidth}{!}{%
		\begin{tabular}{lc}
			\hline
			\textbf{Concepto}        & \textbf{Coste} \\ \hline
			Internet      & 150,00€      \\
			Electricidad    & 175,00€        \\ \hline
			\textbf{Total}   & 325,00€      \\ \hline
		\end{tabular}%
	}
\end{table}

A partir de estos costes obtenemos el coste total del proyecto:
\begin{table}[H]
	\centering
	\caption{Costes totales}
	\resizebox{0.5\textwidth}{!}{%
		\begin{tabular}{lc}
			\hline
			\textbf{Concepto}        & \textbf{Coste} \\ \hline
			Personal      & 6.697,90€      \\
			Hardware    & 45,83€           \\
			Varios     & 325,00€      \\ \hline
			\textbf{Total 5 meses}   & 7.068,73€      \\ \hline
		\end{tabular}%
	}
\end{table}

Para hacer rentable el desarrollo de la aplicación teniendo en cuenta que es de código abierto se podría adoptar un modelo SaaS (Software as a Service) en el que los clientes paguen una suscripción al dueño del software y este a su vez se encargue del hosting y mantenimiento de la aplicación.
Se podrían ofrecer subscripciones a varios niveles ofreciendo distintos niveles de funcionalidad y servicios a distintos precios.
\subsection{Viabilidad legal}
\subsubsection{Licencias de software}
Para determinar la licencia software que va a utilizar la aplicación hay que tener en cuenta las licencias utilizadas por las dependencias que utiliza la aplicación.

\begin{table}[H]
	\caption{Licencias del software utilizado}
	\resizebox{\textwidth}{!}{%
		\begin{tabular}{|c|c|c|}
			\hline
			\textbf{Software}   & \textbf{Descripción}                                                                                                         & \textbf{Licencia} \\ \hline
			Spring Framework    & Framework para aplicaciones web                                                                                              & Apache 2.0        \\ \hline
			Tomcat Embed Jasper & \begin{tabular}[c]{@{}c@{}}Implementación de Tomcat que incluye\\ Jasper, el parser de JSP de Tomcat\end{tabular}            & Apache 2.0        \\ \hline
			JUnit               & Framework para tests unitarios en Java                                                                                       & EPL               \\ \hline
			Apache Commons IO   & \begin{tabular}[c]{@{}c@{}}Librería de utilidades varias (usado en \\ traducción de imágenes a arrays de bytes)\end{tabular} & Apache 2.0        \\ \hline
			Apache Log4j        & Librería para registro de logs                                                                                               & Apache 2.0        \\ \hline
			Bootstrap           & \begin{tabular}[c]{@{}c@{}}Librerías CSS y JavaScript para \\ páginas web\end{tabular}                                       & MIT               \\ \hline
			Plotly.js           & \begin{tabular}[c]{@{}c@{}}Librería JavaScript de \\ generación de gráficos\end{tabular}                                     & MIT               \\ \hline
		\end{tabular}%
	}
\end{table}
La licencia pública de Eclipse (EPL) es compatible con estas licencias, a continuación se mencionan sus posibilidades y obligaciones.
\begin{itemize}
	\item\textbf{Permite:} uso, reproducción, distribución, modificación, uso comercial y uso de patentes.
	\item\textbf{Obliga a:} revelar la fuente y el autor, mantener la misma licencia al redistribuir el software, distribuir el software libre de regalías.
	\item\textbf{No permite:} responsabilizar al autor o contribuidores por posibles daños, utilizar marcas propiedad del autor para promoción o publicidad.
\end{itemize}
\subsubsection{Uso del nombre ``Moodle'' en la aplicación}
Debido a que Moodle es una marca registrada y no soy un partner registrado de Moodle, tengo que atenerme a una serie de restricciones con respecto al uso de la palabra ``Moodle'' \cite{moodletrademark-2022}:
\begin{itemize}
	\item No puedo usar logos de ``Moodle'' sin consentimiento escrito de Moodle.
	\item No puedo usar ``Moodle'' en el nombre de mi software, ni en el de mi dominio ni en el de mi empresa.
	\item No puedo usar ``Moodle'' en palabras clave relacionadas con la publicidad.
	\item No puedo usar ``Moodle'' para describir servicios alrededor de Moodle de forma que la gente piense que estoy asociado a Moodle cuando no es así.
\end{itemize}

Al enterarme de esto a mitad del desarrollo he tenido que cambiar a ``eLearningQA'' el nombre del proyecto, el del repositorio, el del despliegue en Heroku para que el dominio deje de contener ``Moodle'', y he aclarado en el readme del proyecto y en la memoria que no estoy asociado con Moodle.

\apendice{Especificación de Requisitos}

\section{Introducción}

\section{Objetivos generales}

\section{Catalogo de requisitos}
La lista de requisitos funcionales es la siguiente:
\begin{itemize}
	\item R-01: la aplicación permitirá registrarse al docente para acceder a la información sobre los cursos.
	\item R-02: la aplicación permitirá al profesor desconectarse.
	\item R-03: la aplicación mostrará un listado de los cursos en los que se encuentre el profesor.
	\item R-04: se deben poder generar informes que evalúen la calidad de los cursos en los que se encuentra el profesor.
	\item R-05: se deben poder generar informes globales que resuman los informes de todos los cursos en los que se encuentre el profesor.
\end{itemize}

\section{Especificación de requisitos}
A partir de los requisitos se obtiene el siguiente diagrama de casos de uso:
\imagen{DiagramaCasosUso.PNG}{Diagrama de casos de uso}



\apendice{Especificación de diseño}

\section{Introducción}

\section{Diseño de datos}
En la siguiente tabla se recogen las entidades asociadas a cada una de las consultas para determinar parte de las clases necesarias para el proyecto.
\begin{table}[H]
	\resizebox{\textwidth}{!}{%
		\begin{tabular}{|l|l|}
			\hline
			\textbf{Consulta}                                                                                                                                        & \textbf{Entidades}        \\ \hline
			Las opciones de progreso del estudiante están activadas                                                                                         & Curso            \\ \hline
			Se proporcionan contenidos en diferentes formatos                                                                                               & Curso, Recurso            \\ \hline
			El curso tiene grupos                                                                                                                           & Grupo, Curso     \\ \hline
			El curso tiene actividades grupales                                                                                                             & Actividad, Curso \\ \hline
			\begin{tabular}[c]{@{}l@{}}Los estudiantes pueden ver las condiciones\\ necesarias para completar una actividad\end{tabular}                    & Curso            \\ \hline
			Todas las actividades tienen la misma nota máxima en el calificador                                                                             & Actividad, Curso            \\ \hline
			Los recursos están actualizados                                                                                                                 & Curso, Recurso        \\ \hline
			Fechas de apertura y cierre de tareas son correctas                                                                                             & Actividad, Curso \\ \hline
			Se detallan los criterios de evaluación (rúbricas, ejemplos)                                                                                    & Curso            \\ \hline
			El calificador no tiene demasiado anidamiento                                                                                                   & Curso            \\ \hline
			Los alumnos están divididos en grupos                                                                                                           & Alumno, Grupo    \\ \hline
			\begin{tabular}[c]{@{}l@{}}El profesor responde en los foros dentro del límite de\\ 48 horas lectivas desde que se plantea la duda\end{tabular} & Curso, Foro          \\ \hline
			Se ofrece retroalimentación de las tareas                                                                                                       & Tarea, Curso     \\ \hline
			Las tareas están calificadas                                                                                                                    & Tarea, Curso     \\ \hline
			El calificador muestra cómo ponderan las diferentes tareas                                                                                      & Curso            \\ \hline
			La mayoría de alumnos responden a los feedbacks                                                                                                 & Feedback, Curso  \\ \hline
			Se utilizan encuestas de opinión                                                                                                                & Survey, Curso    \\ \hline
		\end{tabular}%
	}
\end{table}
La siguiente imagen muestra las relaciones entre entidades que utiliza la aplicación.
\imagen{ModeloDatos.png}{Modelo de datos base}
\section{Diseño procedimental}
\imagen{BocetoInterfaz.png}{Boceto de la interfaz}
\section{Diseño arquitectónico}



\apendice{Documentación técnica de programación}

\section{Introducción}
En este apéndice se describe una serie de detalles sobre la programación de la aplicación para que sirvan de ayuda a futuros programadores que utilicen el proyecto. Contiene la estructura de directorios del proyecto en la que se explica el contenido de cada uno, un manual del programador que explica ciertos aspectos que puedan ser de ayuda, los detalles sobre la compilación, instalación, y ejecución en un entorno de desarrollo integrado, y por último, qué pruebas del sistema se han realizado.
\section{Estructura de directorios}
El repositorio está organizado de la siguiente manera:
\begin{itemize}
	\item \textbf{/}: Contiene el README y un .gitignore que evita que se suban los archivos de configuración del IDE al repositorio.
	\item \textbf{/.github/workflows}: Contiene el archivo que establece las tareas de integración continua.
	Las tareas son: instalar Java en la maquina que realiza las tareas, almacenar paquetes de SonarCloud en la caché, almacenar paquetes de Maven en la caché, montar (compilar, ejecutar tests y empaquetar) con maven y analizar con SonarCloud.
	\item \textbf{/Memo/plantillaLatex-master}: Contiene la memoria y los anexos.
	\item \textbf{/Memo/plantillaLatex-master/img}: Contiene las imágenes a usar en memoria y anexos.
	\item \textbf{/Memo/plantillaLatex-master/tex}: Contiene los apartados en formato \LaTeX usados de la memoria y anexos.
	\item \textbf{/Project/prototipo}: Contiene el proyecto de la aplicación web. Incluye el Procfile, un archivo usado por Heroku para el despliegue continuo que indica el comando a ejecutar para arrancar la aplicación.
	\item \textbf{/Project/prototipo/configurations}: Contiene los archivos de propiedades.
	\item \textbf{/Project/prototipo/src}: Contiene el código fuente de la aplicación.
	\item \textbf{/Project/prototipo/src/main}: Contiene el código, los test, y los recursos que utiliza la aplicación.
	\item \textbf{/Project/prototipo/src/main/java/es/ubu/lsi}: Contiene las clases necesarias para el funcionamiento de la aplicación.
	\item \textbf{/Project/prototipo/src/main/java/es/ubu/lsi/model}: Contiene clases Java Bean para poder deserializar los datos extraídos por medio de servicios web.
	\item \textbf{/Project/prototipo/src/main/resources}: Contiene el archivo de propiedades de la aplicación web y otros recursos.
	\item \textbf{/Project/prototipo/src/main/resources/images}: Contiene las imágenes a mostrar en la aplicación web.
	\item \textbf{/Project/prototipo/src/main/resources/json}: Contiene información de cursos en formato json para usarla en los tests.
	\item \textbf{/Project/prototipo/src/main/resources/json/informe}: Contiene los textos estáticos de la tabla del informe de fases en formato json.
	\item \textbf{/Project/prototipo/src/main/resources/static/js}: Contiene las librerías JavaScript que utiliza la aplicación.
	\item \textbf{/Project/prototipo/src/main/webapp/WEB-INF/jsp}: Contiene las páginas de la aplicación en formato jsp.
	\item \textbf{/Project/prototipo/src/test/java/es/ubu/lsi}: Contiene los test sobre las consultas que realiza la aplicación en los cursos.
	\item \textbf{/Project/prototipo/target}: Es el directorio donde Maven deposita los resultados, como las clases compiladas y la aplicación compilada en formato WAR.
	
\end{itemize}
\section{Manual del programador}
\subsection{Entorno}
Para trabajar en el proyecto se necesita:
\begin{itemize}
	\item Java JDK 8
	\item IntelliJ IDEA
\end{itemize}

JDK 8 se puede descargar desde \url{https://www.oracle.com/java/technologies/downloads/#java8} si se está en posesión de una cuenta de Oracle.

La versión gratuita de IDEA se puede descargar desde \url{https://www.jetbrains.com/es-es/idea/download}.

\subsection{Obtención del código fuente}
Primero se debería crear un \textit{fork} del repositorio principal en \url{https://github.com/RobertoArastiBlanco/eLearningQA}.
IDEA tiene Git integrado, lo que facilita acciones como los \textit{pull}, \textit{push} y \textit{commit}, pero se puede usar la herramienta de control de versiones que se prefiera para clonar el \textit{fork} y así tenerlo descargado en nuestro equipo.
\imagen{IDEAGit.png}{Desplegable de Git en IDEA}
Para trabajar en el proyecto se debe abrir desde IDEA utilizando \texttt{/Project/prototipo} como directorio raíz.
\subsection{Propósito de las clases principales}
\begin{itemize}
	\item \textbf{Application:} es la clase principal de la aplicación, arranca la aplicación web.
	\item \textbf{SpringController:} indica mediante anotaciones y funciones qué devolver cuando el usuario accede a cada página. En este caso solo devuelve páginas jsp e imágenes, pero se podrían devolver otros tipos de recurso, como archivos de audio o documentos.
	\item \textbf{ELearningQAFacade:} hace de intermediario entre el usuario y la clase WebServiceClient.También contiene funciones para generar tablas y sus contenidos a partir de la información obtenida. 
	\item \textbf{WebServiceClient:} contiene el código de las comprobaciones y realiza las llamadas REST al servidor de Moodle para obtener la información necesaria.
	\item \textbf{FacadeConfig:} almacena la configuración elegida por el usuario junto con el host de Moodle y se guarda como atributo de la fachada para poder pasarlo a las funciones de WebServiceClient y que pueda cambiar su comportamiento según la configuración.
	\item \textbf{AlertLog:} genera y almacena mensajes de alerta, se pasa como parámetro a las funciones de WebServiceClient para que puedan almacenar los mensajes para la lista de mejoras del informe.
	\item \textbf{RegistryIO:} se encarga de almacenar y acceder a la información guardada en csv utilizada para generar los gráficos. También contiene una función que genera un gráfico a partir de los datos de un csv.
	\item \textbf{AnalysisSnapshot:} la clase que representa el conjunto de datos almacenado en cada línea de los archivos csv, facilita el trabajo de almacenamiento de datos y generación de gráficos.
\end{itemize}
\subsection{Integración continua}
El proyecto contiene un archivo llamado \texttt{maven.yml} en \texttt{/.github/workflows} en el que se indica los pasos a seguir tras realizar un \textit{push} al repositorio. Vamos a explicarlos brevemente. \imagen{MavenYML.png}{Trabajos de integración continua}
Lo primero que se puede apreciar es que la máquina virtual que se encarga de la integración continua utiliza la versión más reciente de Ubuntu. Después, se instala JDK 11 (antes se utilizaba JDK 8 para la integración continua, pero tras añadir SonarCloud al proceso, JDK 11 pasó a ser necesario), luego se almacenan en la memoria caché de la máquina virtual los paquetes de SonarCloud y Maven. Finalmente, se construye y analiza el proyecto, se indica la ubicación del archivo \texttt{pom.xml} debido a que el directorio raíz del proyecto no coincide con el del repositorio (después veremos otro problema causado por esto).
\subsection{Herramienta de calidad de código}
Para analizar la calidad del código del proyecto hemos utilizado SonarCloud.
Para integrar SonarCloud en un proyecto se puede hacer desde \url{https://sonarcloud.io/projects/create} y se hace relativamente fácil, ya que durante el proceso se te guía paso a paso.

La página principal resume el estado del proyecto respecto a errores, \textit{code smells}, seguridad, y cantidad de código duplicado. También muestra el estado del proyecto según el \textit{Quality Gate}, que indica si se cumple una serie de condiciones establecidas.
\imagen{SonarCloud.png}{Página del proyecto en SonarCloud en \url{https://sonarcloud.io/summary/overall?id=RobertoArastiBlanco_MoodleQA}}
Es muy útil a la hora de arreglar los errores encontrados, ya que posee variedad de filtros para los errores, explica el motivo del error, posible causa, e incluso suele sugerir formas de solventarlo.
\imagen{SonarCloudIssues.png}{Apartado de problemas de código}

\section{Compilación, instalación y ejecución del proyecto}
\subsection{Ejecución en local}
Para ejecutar la aplicación en local solo es necesario ejecutar la clase principal: Application, IDEA se encarga de compilar las clases necesarias de forma automática. Con la aplicación en ejecución, se puede acceder a ella en \url{http://localhost:8080/}.
\imagen{EjecucionLocal.png}{Ejecución de la clase principal}
\imagen{EjecucionLocalhost.png}{Ejecución en local}
\subsection{Despliegue en Heroku}
A la hora de desplegar la aplicación en Heroku, existen algunas consideraciones a tener en cuenta para conseguirlo.

Para indicar a Heroku cómo ejecutar la aplicación, el directorio raíz del proyecto debe contener un archivo llamado \texttt{Procfile} (sin extensión de archivo) con el comando ``web: java -jar -Dserver.port=\$PORT target/prototipo-0.4-SNAPSHOT.war'' (el nombre del ejecutable puede cambiar al cambiar la versión de la aplicación en \texttt{pom.xml})

Debido a que el directorio raíz de la aplicación difiere del directorio raíz del repositorio, es necesario el uso de un \textit{buildpack} específico, un script que se ejecuta en el despliegue y configura el entorno de ejecución.

Se debe especificar la ruta del directorio raíz de la aplicación asociada a la clave \texttt{PROJECT\_PATH} en variables de configuración y añadir los buildpacks \texttt{\url{https://github.com/timanovsky/subdir-heroku-buildpack.git}} y \texttt{heroku/java}
\imagen{Buildpack.png}{Configuración necesaria para el despliegue}

En la pestaña ``Deployment'' debes elegir GitHub como método de despliegue. Al indicar el repositorio, utiliza el \textit{fork} creado, porque se te pedirá registrarte con GitHub para obtener un \textit{token OAuth} para tener acceso al repositorio y se necesita ser colaborador para tener acceso.

\imagen{Deployment.png}{Configuración del despliegue automático}

\section{Pruebas del sistema}
Los test se realizan con JUnit.
La clase ELearningQAFacadeTest contiene los test unitarios de la clase ELearningQAFacade. Al principio se realizaban los test contra la página school.moodledemo.net, accesible por cualquiera para realizar pruebas en un entorno Moodle, el problema de hacer esto es que ralentizaba la ejecución de los test notablemente y el resultado de estos estaba sujeto a cambios que pudiera hacer otra persona usando la aplicación para hacer pruebas.
La solución que se encontró a este problema fue refactorizar los test para que extrajeran los datos de los cursos de prueba de archivos JSON almacenados en un directorio de recursos de la aplicación. Para generar estos archivos creé de forma temporal una función llamada ``guardarComoJSON'' que convertía un objeto cualquiera a JSON y lo almacenaba en el directorio ``resources/json''.
La clase ELearningQAFacadeTest contiene una función con la anotación ``@BeforeAll'' que indica que debe ejecutarse antes que el resto de funciones de test. La función en cuestión inicializa la fachada y carga el conjunto de datos de cursos que se encontraba almacenado en formato JSON.
\apendice{Documentación de usuario}

\section{Introducción}
En este apéndice se incluyen los detalles de la aplicación que conciernen al usuario, como los requisitos para utilizar la aplicación y un manual de usuario que también existe en formato web y puede ser accedido desde la página de login de la aplicación.
\section{Requisitos de usuarios}
Para utilizar la aplicación como usuario se necesita:
\begin{itemize}
	\item Una cuenta en un servidor de Moodle con rol de profesor en los cursos a consultar
	\item Java 8 si se ejecuta la aplicación en local
	\item Conexión a internet para acceder al despliegue de la aplicación o al servidor de Moodle correspondiente si se ejecuta en local
\end{itemize}
\section{Instalación}

\section{Manual del usuario}

\subsubsection{Pantallas:}

\textbf{Página de login:} la página de login contiene un formulario de
inicio de sesión, un link ''Acerca de'', un link ''Manual de usuario'' que
muestra la versión web de este manual, y un link de contacto.
\imagen{Login.png}{Página de login}

\textbf{Página principal:~}la página principal contiene una cabecera que
indica el nombre de usuario con el que se ha registrado, un botón para
desconectarse, un enlace para generar un informe global de todos los
cursos, y una tabla con una lista de enlaces para generar los informes
específicos de cada curso con una caja de búsqueda.
\imagen{ListaCursos.png}{Página principal}

\textbf{Página de informe:} la página de informe genera el informe
dependiendo de cómo se haya accedido a esta, ya sea mediante el enlace
''Generar informe global'' o por medio de los enlaces a informes de cursos
específicos.~

Los informes contienen los siguientes elementos:

\begin{itemize}
	\item
	\textbf{Puntuaciones por rol y perspectiva:} el informe contiene una
	matriz de roles y perspectivas que da distintas puntuaciones a cada
	uno de los roles (Diseñador, Facilitador, y Proveedor) en cada una de
	las tres perspectivas definidas (Pedagógica, Tecnológica, y
	Estratégica) según los resultados del análisis.
	\item
	\textbf{Gráfico de evolución:} la aplicación almacena estas puntuaciones cada vez que se carga un informe para generar un gráfico interactivo de la evolución de estos valores.
	\imagen{Evolucion.png}{Página de evolución del rendimiento}
	\item
	\textbf{Informe de fases:} el informe contiene también una tabla que
	muestra los resultados en cada una de las comprobaciones del análisis y
	las agrupa en las distintas fases del diseño instruccional (Diseño,
	Implementación, Realización, y Evaluación) además de dar una
	puntuación general.
	\item
	\textbf{Lista de mejoras:} por último, el informe muestra una lista de
	motivos concretos que producen resultados negativos en el análisis
	para subsanarlos o tenerlos en cuenta de cara al futuro. Haciendo clic en las filas de la tabla se puede filtrar la lista de mejoras.
	\imagen{InformeFases.png}{Página de informe de fases}
\end{itemize}

El informe especifico muestra los resultados de las comprobaciones en el
informe de fases de forma absoluta, es decir, que solo indica si el
criterio se cumple o no, sin embargo, en el informe global, se muestra
un resultado u otro ~dependiendo de la cantidad de cursos que satisfacen
dicho criterio.


\subsubsection{Acciones del usuario:}

\begin{itemize}
	\item
	\textbf{Login:} para acceder a la aplicación son necesarias las
	credenciales de acceso a una cuenta de la plataforma Moodle a la que
	accede la aplicación (en el caso del prototipo es Mount ~Orange
	School). Debe introducir su usuario y contraseña en los campos
	''Username'' y ''Password''. Si quiere vaciar los campos pulse el botón
	''Borrar''. Para acceder a la página principal, pulse el botón ''Entrar''
	tras haber introducido sus credenciales.
	\item
	\textbf{Desconectar:} desde la página principal, si desea finalizar su
	sesión, pulse el botón ''Desconectar''. Esto invalidará sus credenciales
	y le impedirá acceder a la aplicación hasta que se registre de nuevo
	con unas credenciales válidas.
	\item
	\textbf{Generar informe específico:} la página principal muestra una
	tabla con todos los cursos en los que se encuentra matriculado el
	usuario registrado en formato de enlace. Al clicar un enlace, se
	generará un informe en una pestaña aparte del navegador que mostrará
	los resultados del análisis que ha realizado la aplicación sobre el
	curso correspondiente.
	\item
	\textbf{Generar informe global:} en la página principal hay un enlace
	llamado ''Generar informe global''. Al hacer clic sobre este, se
	generará un informe en una pestaña aparte del navegador que mostrará
	un resumen de los análisis de todos los cursos en los que se encuentra
	matriculado el profesor.
\end{itemize}

\subsubsection{Explicación de las comprobaciones de los informes:}

Las siguientes comprobaciones están relacionadas con los roles, fases, y
perspectivas mencionados anteriormente. Los distintos procesos del
diseño instruccional se encuentran divididos en fases, ~con ciertas
perspectivas en mente, y son responsabilidad directa o indirecta de
ciertos roles. Al estar las comprobaciones ligadas a esos procesos se
muestran agrupadas por fases, y después de la explicación se indican los
roles responsables e involucrados, además de las perspectivas
correspondientes.

\textbf{Diseño:}

\begin{itemize}
	\item
	\textbf{Las opciones de progreso del estudiante están activadas}: se
	comprueba que estén habilitadas las opciones de progreso de los
	estudiantes en el curso.~{Responsable:} Diseñador
	{Involucrados:} Facilitador {Perspectivas:} Pedagógica
	\item
	\textbf{Se proporcionan contenidos en diferentes formatos:} se
	comprueba que haya variedad de formatos en los recursos del curso.
	{Responsable:} Diseñador {Involucrados:} Facilitador y
	Proveedor {Perspectivas:} Pedagógica y Tecnológica
	\item
	\textbf{El curso tiene grupos:} se comprueba que existan grupos
	definidos en el curso. {Responsable:} Diseñador
	{Involucrados:} Facilitador y Proveedor {Perspectivas:}
	Pedagógica
	\item
	\textbf{El curso tiene actividades grupales:} se comprueba que existan
	actividades con entrega grupal habilitada en el curso.
	{Responsable:} Diseñador {Involucrados:} Facilitador y
	Proveedor {Perspectivas:} Pedagógica
	\item
	\textbf{Los estudiantes pueden ver las condiciones necesarias para
		completar una actividad:} se comprueba que esté habilitada la opción
	de mostrar las condiciones para completar una actividad en el curso.
	{Responsable:} Diseñador {Involucrados:} Facilitador y
	Proveedor {Perspectivas:} Pedagógica
	\item
	\textbf{Todas las actividades tienen la misma nota máxima en el
		calificador:} se comprueba que exista una consistencia en las notas
	máximas de los items de calificación (tareas, entregas, cuestionarios)
	del curso. {Responsable:} Diseñador {Involucrados:}
	Facilitador y Proveedor {Perspectivas:} Pedagógica
\end{itemize}

\textbf{Implementación:}

\begin{itemize}
	\item
	\textbf{Los recursos están actualizados:} se comprueba que los
	recursos del curso tengan una fecha de creación reciente.
	{Responsable:} Diseñador {Involucrados:} Facilitador y
	Proveedor {Perspectivas:} Pedagógica y Tecnológica
	\item
	\textbf{Fechas de apertura y cierre de tareas son correctas:} se
	comprueba que las fechas de apertura y cierre de tareas y
	cuestionarios no se solapen de forma erronea con las fechas de inicio
	y fin del curso. {Responsable:} Facilitador
	{Involucrados:} Diseñador y Proveedor {Perspectivas:}
	Pedagógica y Tecnológica
	\item
	\textbf{Se detallan los criterios de evaluación:} se comprueba que
	exista en al menos una actividad una rúbrica o una guía de calificación
	en el curso.~{Responsable:} Diseñador {Involucrados:}
	Facilitador y Proveedor {Perspectivas:} Pedagógica y Tecnológica
	\item
	\textbf{El calificador no tiene demasiado anidamiento:} se comprueba
	que la estructura de las categorías del calificador no sea demasiado
	enrevesada. {Responsable:} Diseñador {Involucrados:}
	Facilitador y Proveedor {Perspectivas:} Pedagógica y Estratégica
	\item
	\textbf{Todos los alumnos están en algún grupo:} se comprueba que cada
	alumno pertenezca a un grupo. {Responsable:} Proveedor
	{Involucrados:} Diseñador {Perspectivas:} Tecnológica y
	Estratégica
\end{itemize}

\textbf{Realización:}

\begin{itemize}
	\item
	\textbf{El profesor responde en los foros dentro del límite de 48
		horas lectivas desde que se plantea la duda:} se comprueba que no
	hayan preguntas por parte de alumnos que estén sin responder en un
	tiempo razonable. {Responsable:} Facilitador
	{Involucrados:} Diseñador y Proveedor {Perspectivas:}
	Pedagógica y Tecnológica
	\item
	\textbf{Se ofrece retroalimentación de las tareas:} se comprueba que
	el profesor deje comentarios en la mayoría de calificaciones que haga.
	{Responsable:} Facilitador {Involucrados:} Diseñador y
	Proveedor {Perspectivas:} Pedagógica y Tecnológica
	\item
	\textbf{Las tareas están calificadas:} se comprueba que no hayan
	entregas de alumnos que hayan pasado una semana sin calificación.
	{Responsable:} Facilitador {Involucrados:} Diseñador y
	Proveedor {Perspectivas:} Pedagógica y Tecnológica
	\item
	\textbf{El calificador muestra cómo ponderan las diferentes tareas:}
	se comprueba que el calificador muestre los pesos de los items de
	calificación. {Responsable:} Facilitador {Involucrados:}
	Diseñador y Proveedor {Perspectivas:} Pedagógica y Tecnológica
\end{itemize}

\textbf{Evaluación:}

\begin{itemize}
	\item
	\textbf{La mayoría de alumnos responden a los feedbacks:} se comprueba
	que no hayan muchos alumnos que no respondan a los feedbacks.
	{Responsable:} Proveedor {Involucrados:} Diseñador y
	Facilitador {Perspectivas:} Pedagógica, Tecnológica, y
	Estratégica
	\item
	\textbf{Se utilizan encuestas de opinión:} se comprueba que el curso
	contenga encuestas de opinión. {Responsable:} Proveedor
	{Involucrados:} Diseñador y Facilitador {Perspectivas:}
	Pedagógica, Tecnológica, y Estratégica
\end{itemize}




\bibliographystyle{plain}
\bibliography{bibliografiaAnexos}

\end{document}
